\section{Síntesis y ampliación}

La conclusión central de esta clase puede resumirse en una frase: la capacidad de planning de Codex no es un botón único, sino un conjunto de mecanismos por capas. Plan Mode se encarga de aclarar la especificación antes de actuar; \texttt{update\_plan} se encarga de mantener el live progress durante la ejecución; y el todo en archivo se encarga de desacoplar el estado de largo plazo de la context window.

Al final, quien dio la clase señala que, a medida que la capacidad de los modelos aumenta, una herramienta tan minimalista como \texttt{update\_plan} ya basta para que el modelo lleve a cabo creación, avance de estado y replan complejos. Lo que en verdad importa no es el número de herramientas, sino si el schema de la herramienta está alineado con la capacidad semántica del modelo.

\begin{importantbox}{Takeaway final}
Un buen agent harness no necesariamente tiene que acumular muchas herramientas especializadas. Una herramienta lo bastante general, con semántica clara y estado visible, combinada con un modelo fuerte, suele cubrir un espacio de conducta mucho mayor. El \texttt{update\_plan} de Codex es el representante de este tipo de diseño.
\end{importantbox}

Para mantener nuestro repositorio de notas de curso también podemos tomar esta idea como referencia: para tareas cortas, avanzar con un plan en el contexto; para tareas largas, fijar el estado con un todo en archivo; y cada vez que nos encontremos con un bloqueo de descarga, subtítulos, transcripción, compilación, etc., no quedarnos solo en la memoria oral, sino escribirlo de vuelta en la cola, para formar un flujo de trabajo recuperable.

\end{document}
